\documentclass[a4paper, serif, 10pt]{moderncv}

%% moderncv
% style and color
\moderncvstyle{banking}
\moderncvcolor{black}

%% page layout/margins
\usepackage[top=2.5cm, bottom=2.5cm, left=1.5cm, right=1.5cm]{geometry}

\nopagenumbers{}

%% Babel config for Dutch language
% ref:
% - https://latex3.github.io/babel/guides/locale-dutch.html
\usepackage[dutch]{babel}

%% fontspec
\usepackage{fontspec}

\IfFontExistsTF{Linux Libertine}{
  \setmainfont[Ligatures={TeX, Common}, Numbers=OldStyle]{Linux Libertine}
}{}
\IfFontExistsTF{Linux Libertine O}{
  \setmainfont[Ligatures={TeX, Common}, Numbers=OldStyle]{Linux Libertine O}
}{}

%% CTeX
% CJK support, used to show CJK characters in the resume
%
% - fontset=none: disable builtin fontset but instead set the CJK font manually
% - heading=false: disable ctex heading
% - punct=kaiming: use kaiming punctuations styles for CJK
% - scheme=plain: use plain scheme, do not override \normalsize font size
% - space=auto: space settings for CJK characters
%
% ref:
% - http://ctan.mirrorcatalogs.com/language/chinese/ctex/ctex.pdf
\usepackage[UTF8, heading=false, punct=kaiming, scheme=plain, space=auto]{ctex}

\IfFontExistsTF{Noto Serif CJK SC}{
  \setCJKmainfont{Noto Serif CJK SC}
}{}
\IfFontExistsTF{Noto Sans CJK SC}{
  \setCJKsansfont{Noto Sans CJK SC}
}{}

%% line spacing
\usepackage{setspace}
\setstretch{1.125}

%% URL styling - use normal text instead of monospace
\urlstyle{same}

% Auto-underline all links
% Defer until \begin{document} so hyperref is already loaded
\AtBeginDocument{%
  \let\oldhref\href
  \renewcommand{\href}[2]{\oldhref{#1}{\underline{#2}}}%
}

%% Patch \httplink and \httpslink to handle full URLs
% The original moderncv macros always prepend http:// or https:// to URLs.
% This causes issues when the URL already has a protocol (e.g., https://example.com)
% resulting in malformed URLs like https://https://example.com.
% This patch checks if the URL already contains "://" and uses it directly if so.
% Note: We use \str_set:Nx (x-expansion) to expand macros like \@homepage first.
\ExplSyntaxOn
\str_new:N \g_yamlresume_url_str

\RenewDocumentCommand{\httpslink}{O{}m}{
  \str_set:Nx \g_yamlresume_url_str {#2}
  \str_if_in:NnTF \g_yamlresume_url_str {://}
    { \url{#2} }
    { \url{https://#2} }
}

\RenewDocumentCommand{\httplink}{O{}m}{
  \str_set:Nx \g_yamlresume_url_str {#2}
  \str_if_in:NnTF \g_yamlresume_url_str {://}
    { \url{#2} }
    { \url{http://#2} }
}
\ExplSyntaxOff

\name{Lars van den Berg}{}
\title{Software Engineer met passie voor schaalbare cloud-native systemen}
\phone[mobile]{+31 6 12345678}
\email{lars@vandenberg.dev}
\homepage{https://www.larsvandenberg.dev}

\address{Oudegracht 123, Utrecht, Utrecht, Nederland, 3511 AB}{}{}

\extrainfo{{\small \faGithub}\ \href{https://github.com/larsvandenberg}{@larsvandenberg} {} {} {} • {} {} {}
{\small \faLinkedin}\ \href{https://www.linkedin.com/in/larsvandenberg}{@larsvandenberg}}

\begin{document}

\maketitle

\section{Basis Informatie}

\cvline{}{Software Engineer met ruim 8 jaar ervaring in het ontwerpen en bouwen van
schaalbare, cloud-native applicaties. Specialist in backend-ontwikkeling met
Python, Java en Go, en uitgebreide ervaring met AWS en Kubernetes.
%
\begin{itemize}
\item Ontwerpen van event-gedreven microservices met hoge beschikbaarheid
\item Automatiseren van deployment pipelines met CI/CD en Infrastructure as Code
\item Sterke focus op codekwaliteit, testbaarheid en maintainability
\item Ervaren in het begeleiden van junior ontwikkelaars en agile samenwerking
\end{itemize}}
\section{Opleidingen}

\cventry{sep 2013–aug 2015}
        {Master, Informatica, Score: 8.0}
        {Technische Universiteit Delft}
        {\url{https://www.tudelft.nl}}
        {}
        {Afstudeeronderzoek gericht op schaalbare, event-gebaseerde architecturen.
Ontwikkelde een prototype voor een event-sourcing platform dat gebruikt
werd door meerdere onderzoeksgroepen.
\textbf{Cursussen}: Distributed Systems, Machine Learning, Software Architecture, Big Data Processing}

\cventry{sep 2010–aug 2013}
        {Bachelor, Informatica, Score: 7.8}
        {Hogeschool Utrecht}
        {\url{https://www.hu.nl}}
        {}
        {Bachelorstage uitgevoerd bij een fintech-bedrijf, waar ik een
realtime-transactiemonitor bouwde met Java en WebSocket.
\textbf{Cursussen}: Objectgeoriënteerd programmeren, Databases, Software Engineering, Webtechnologie}
\section{Werkervaring}

\cventry{feb 2022–Heden}
        {Senior Backend Engineer}
        {CloudNimbus B.V.}
        {\url{https://www.cloudnimbus.nl}}
        {}
        {\begin{itemize}
\item Verantwoordelijk voor de architectuur en ontwikkeling van meerdere event-gedreven microservices
\item Migreerde een monolithische applicatie naar Kubernetes en AWS EKS
\item Introduceerde Infrastructure as Code met Terraform, wat de deploymenttijd met 70\% verminderde
\item Begeleidde een team van vijf ontwikkelaars middels code reviews en technische workshops
\item Werkt nauw samen met productowners om technische roadmap en sprints te bepalen
\end{itemize}
\textbf{Trefwoorden}: Microservices, Kubernetes, AWS, Terraform, Java, Mentoring}

\cventry{jun 2018–jan 2022}
        {Software Engineer}
        {DataVault Solutions}
        {\url{https://www.datavault.nl}}
        {}
        {\begin{itemize}
\item Ontwikkelde en onderhield een API-platform voor veilige data-uitwisseling
\item Implementeerde OAuth2 en JWT-authenticatie voor meerdere externe partijen
\item Optimaliseerde databasequery's, wat leidde tot een significante verlaging van de responstijden
\item Droeg bij aan de inrichting van CI/CD-pipelines met GitLab CI
\item Werkte in een multidisciplinair team volgens de Scrum-methodiek
\end{itemize}
\textbf{Trefwoorden}: Python, FastAPI, PostgreSQL, OAuth2, GitLab CI}

\cventry{jul 2016–mei 2018}
        {Junior Software Engineer}
        {BrightCode}
        {\url{https://www.brightcode.nl}}
        {}
        {\begin{itemize}
\item Ontwikkelde webapplicaties met Java Spring Boot en Angular
\item Vulde user stories aan met technische ontwerpen en testplannen
\item Schreef unit- en integratietests om de kwaliteit van de software te waarborgen
\item Participeerde in dagelijkse stand-ups en sprint retrospectives
\end{itemize}
\textbf{Trefwoorden}: Java, Spring Boot, Angular, JUnit, Scrum}
\section{Talen}

\cvline{Nederlands}{Moedertaal of tweetalige vaardigheid \hfill \textbf{Trefwoorden}: Moedertaal}
\cvline{Engels}{Volledige professionele vaardigheid \hfill \textbf{Trefwoorden}: Cambridge C2, Zakelijk Engels}
\cvline{Duits}{Beperkte werkvaardigheid \hfill \textbf{Trefwoorden}: Conversational}
\section{Vaardigheden}

\cvline{Backend-ontwikkeling}{Expert \hfill \textbf{Trefwoorden}: Python, Java, Go, Node.js, Spring Boot, FastAPI}
\cvline{Cloud \& DevOps}{Geavanceerd \hfill \textbf{Trefwoorden}: AWS, Kubernetes, Docker, Terraform, GitHub Actions, ArgoCD}
\cvline{Databases}{Geavanceerd \hfill \textbf{Trefwoorden}: PostgreSQL, MySQL, MongoDB, Redis, DynamoDB}
\cvline{Frontend-ontwikkeling}{Intermediair \hfill \textbf{Trefwoorden}: React, TypeScript, Tailwind CSS, Next.js}
\section{Onderscheidingen}

\cventry{aug 2015}
        {Technische Universiteit Delft}
        {Afstudeerprijs Informatica}
        {}
        {}
        {Ontvangen voor het afstudeeronderzoek naar schaalbare event-gebaseerde
systemen, inclusief de publicatie van een bijbehorend prototype als
open-source project.}
\section{Certificaten}

\cventry{mrt 2021}
        {Amazon Web Services}
        {AWS Certified Solutions Architect – Associate}
        {\url{https://www.credly.com/badges/aws-solutions-architect}}
        {}
        {}

\cventry{jun 2022}
        {Cloud Native Computing Foundation}
        {Certified Kubernetes Application Developer}
        {\url{https://www.cncf.io/certification/ckad/}}
        {}
        {}
\section{Publicaties}

\cventry{apr 2020}
        {Event-driven microservices in de praktijk}
        {Java Magazine NL}
        {\url{https://www.java-magazine.nl/publicaties/event-driven-microservices}}
        {}
        {Praktijkartikel over het ontwerpen van event-driven microservices, inclusief
lessons learned bij de migratie van een legacy-systeem naar een
cloud-native architectuur.}
\section{Projecten}

\cventry{jan 2023–Heden}
        {Open-source command-line tool voor het beheren van Docker-containers.
Voorziet in overzicht van actieve containers, logweergave en healthchecks.
}
        {ContainerCLI}
        {\url{https://github.com/larsvandenberg/containercli}}
        {}
        {Gebouwd in Go met behulp van Cobra en de Docker Engine API. Het project
heeft inmiddels meer dan 200 GitHub-sterren en wordt actief gebruikt door
ontwikkelaars wereldwijd.
\textbf{Trefwoorden}: Go, Docker, CLI, Cobra}

\cventry{sep 2021–dec 2021}
        {Persoonlijke portfolio-website met projectenoverzicht, blog en
contactformulier.
}
        {Portfolio-website}
        {\url{https://www.larsvandenberg.dev}}
        {}
        {Ontworpen en geïmplementeerd met Next.js, TypeScript en Tailwind CSS.
De website scoort 100 op Lighthouse voor prestaties, toegankelijkheid en
SEO.
\textbf{Trefwoorden}: Next.js, TypeScript, Tailwind CSS, Vercel}
\section{Interesses}

\cvline{Duurzaamheid}{Fietsen, Klimaattechnologie, Zonne-energie}
\cvline{Technologie}{Open source, Self-hosting, Retrocomputing}
\section{Vrijwilligerswerk}

\cventry{jan 2023–Heden}
        {Open-source bijdrager}
        {NL Design System}
        {\url{https://www.nldesignsystem.nl}}
        {}
        {Draag bij aan toegankelijke, herbruikbare UI-componenten voor Nederlandse
overheidswebsites. Geef feedback op code, schrijf documentatie en lever
verbeteringen aan op het gebied van toegankelijkheid.}

\end{document}